% Include in an article with amsmath, amssymb, amsthm and hyperref.
% Companion cross-references refer to paper-unified/Supplement.tex.

\section{The finite lift and the positivity it requires}
\label{sec:insight-statement}

Consider two separated quantum systems, with two binary measurements
on each. Write their eight expectations as
\[
 q=(a_0,a_1,b_0,b_1,E_{00},E_{01},E_{10},E_{11}),\qquad
 p(\alpha,\beta\mid x,y)
   ={1+\alpha a_x+\beta b_y+\alpha\beta E_{xy}\over4}.
\]
Let $\mathcal Q_{222}$ be the set obtained from arbitrary
finite-dimensional bipartite states and binary measurements.

\begin{theorem}\label{thm:insight-finite-lift}
There are finite integers $N,p$ and real symmetric matrices of order $N$,
denoted by $L_0,\ldots,L_8$ and $M_1,\ldots,M_p$, such that
\[
 \mathcal Q_{222}
 =\left\{q\in\mathbb R^8:\ \exists y\in\mathbb R^p,\quad
       L_0+\sum_{\nu=1}^8q_\nu L_\nu+
                         \sum_{j=1}^py_jM_j\succeq0\right\}.
\]
The matrix entries can be chosen to be real algebraic numbers.
\end{theorem}

The decisive step is to preserve positivity when the Bell inequality
and the experiment approach a singular contact independently.
A stationary contact already determines the global angular minimum
through a balanced sum of eight Laurent monomials. To obtain a
finite lift, we must also express the full state quadratic form as
a sum of squares with the two sets of parameters kept separate.
The proof below explains this passage. Its source-specific estimates,
including simultaneous vanishing coefficients and projective
endpoints, are proved in the complete companion.

The two-qubit reduction gives a fixed polynomial model for the set.
Two binary measurements can be dilated using one common embedding,
then decomposed into common blocks of dimension at most two.
In each block the observables can be written
\[
 A_0=B_0=Z,\quad A_1=c_AZ+s_AX,\quad B_1=c_BZ+s_BX,
 \qquad c_A^2+s_A^2=c_B^2+s_B^2=1,
\]
with $X,Z$ the real Pauli matrices. Because these matrices are real,
the real and imaginary parts of a complex state give a convex
mixture of real states. Thus, for the polynomial expectation map
$\Phi$ defined by these observables,
\begin{equation}\label{eq:insight-finite-model}
 \mathcal Q_{222}=\operatorname{conv}\Phi(S^1\times S^1\times S^3)
 =\left\{\sum_{i=1}^{9}w_i\Phi(z_i):
 w_i\ge0,\ \sum_iw_i=1,\ z_i\in S^1\times S^1\times S^3\right\}.
\end{equation}
The first reduction is the usual real-qubit description of this
scenario~\cite{Masanes}; Chapter~2 proves the dilation,
block decomposition and both inclusions, including commuting blocks.
The nine-atom bound is Carath\'eodory's theorem in eight dimensions.
Equation~\eqref{eq:insight-finite-model} makes the set compact and
semialgebraic and gives a finite formula over $\mathbb Q$ for its
extension to any real closed field.

Let $P$ be the cone of affine functions nonnegative on
$\mathcal Q_{222}$. Such a function is the slack of a Bell bound.
For separately labelled real closed extensions $K_c,K_t$ of
$\mathbb R$, put $T=K_c\otimes_{\mathbb R}K_t$ and define
\begin{equation}\label{eq:insight-tensor-slack}
 [f,q]=f_0\otimes1+\sum_{\nu=1}^8f_\nu\otimes q_\nu,
 \qquad f\in P_{K_c},\quad q\in(\mathcal Q_{222})_{K_t}.
\end{equation}
The subscripts denote transfer of the fixed finite formulas.
Scheiderer's Corollary~3.19~\cite{Scheiderer} identifies finite
semidefinite representability with the following condition:
for every real closed extension $K/\mathbb R$, every
$f\in P_K$ and every $q\in(\mathcal Q_{222})_K$, their evaluation
in two labelled copies of $K$ is a finite sum of squares in
$K\otimes_{\mathbb R}K$. We prove the stronger statement
\begin{equation}\label{eq:insight-universal-goal}
 [f,q]\in\Sigma T^2\qquad
 \text{for every }K_c,K_t,\ f\in P_{K_c},\ q\in(\mathcal Q_{222})_{K_t}.
\end{equation}
No algebraic-disjointness assumption is imposed. The number of
squares may depend on these choices; the criterion supplies one
finite lift from the universal assertion.

\section{A stationary contact balances the angular monomials}
\label{sec:insight-angular}

Choose
$S=E_{00}-E_{01}+E_{10}+E_{11}$ and write $H=1-S/2$.
Order its eight losing events as
\begingroup\setlength{\arraycolsep}{3pt}
\[
\begin{array}{c|cccccccc}
\text{weight}&\lambda_{\rm ev}&n&m&C&A&D&\alpha&\beta\\
\text{event}&(++|01)&(+-|00)&(+-|11)&(-+|10)&(-+|00)&(-+|11)&(--|01)&(+-|10).
\end{array}
\]
\endgroup
Every valid slack
which is negative at the corresponding PR vertex has the unique form
\begin{equation}\label{eq:insight-event-source}
 f=dH+4\sum_i c_i p_i,\qquad d>0,\quad c_i\ge0.
\end{equation}
Here $p_i$ are its eight losing-event probabilities; evaluating at
the PR vertex and its eight deterministic neighbours determines
the coefficients. A slack nonnegative at every PR vertex is a
nonnegative combination of event probabilities, whose separated
square certificates are immediate.

For half-angles $u,v$ and a homogeneous real state
$\psi=(W,X,Y,Z)$, clear the positive measurement denominators and
the state norm. With $D_u=1+u^2$, $D_v=1+v^2$, the resulting
quadratic form $F_c=\psi^{\mathsf T}P_c(u,v)\psi$ is
\begin{align}
 F_c={}&4\lambda_{\rm ev}D_u(W+vX)^2+4nD_uD_vX^2
             +4m(vW-X+uvY-uZ)^2\notag\\
 &+4CD_v(uW-Y)^2+4AD_uD_vY^2
             +(4D+2d)(uW+uvX-Y-vZ)^2\notag\\
 &+(4\alpha+2d)D_u(vY-Z)^2
             +(4\beta+2d)D_v(X+uZ)^2-2dD_uD_vZ^2.
 \label{eq:insight-parent}
\end{align}
Chapter~2, Eq.~(RFA10), gives its homogeneous angle
form at all projective endpoints. Each event contributes a
nonnegative amplitude square. The signed CHSH part has just one
negative square, $-2dD_uD_vZ^2$. For $uv\ne0$, the $Z=0$ block
of $P_c$ contains the positive-definite form
\begin{equation}\label{eq:insight-transverse}
 2d\bigl[(uW+uvX-Y)^2+(1+u^2)v^2Y^2+(1+v^2)X^2\bigr].
\end{equation}
Hence $P_c$ has three positive transverse directions and
$P_c\succeq0$ is equivalent to $\det P_c\ge0$.

The determinant has a sparse structure:
\begin{equation}\label{eq:insight-determinant}
 \begin{gathered}
 \det P_c=16(1+u^2)^2(1+v^2)^2u^2v^2L_c(u^2,v^2),\\
 L_c(U,V)=h_c+\sum_{e\in E}q_eU^{e_1}V^{e_2},
 \end{gathered}
\end{equation}
where $E=\{-1,0,1\}^2\setminus\{0\}$ and every $q_e$ is
nonnegative. The eight factored coefficients are Eq.~(SG6) in
Chapter~4. Their signs are the reason a local
stationarity condition determines global physicality.

\begin{lemma}[Balanced angular positivity]\label{lem:insight-balance}
Suppose $q_e,h_c,U_c,V_c\in K_c$, $q_e\ge0$, $U_c,V_c>0$, and
$L_c=L_{c,U}=L_{c,V}=0$ at $(U_c,V_c)$. Then the separated
evaluation of $L_c$ belongs to $\Sigma T^2$ at every positive
$U,V\in K_t$.
\end{lemma}
\begin{proof}
Put $w_e=q_eU_c^{e_1}V_c^{e_2}$. The three contact equations are
\[
 h_c=-\sum_e w_e,\qquad \sum_e w_e e=0.
\]
Decompose this nonnegative balanced vector into minimal positive
dependencies among the eight exponents. A minimal dependency is
an opposite pair, or, up to symmetries of the square, one of
\[
 \{(1,1),(-1,0),(0,-1)\},\qquad
 \{(1,1),(-1,1),(0,-1)\},
\]
with integer weights $(1,1,1)$ and $(1,1,2)$, respectively.
To obtain the decomposition, subtract from any such active
dependency its largest nonnegative multiple. At least one active
coefficient disappears, so this finishes after finitely many steps.

Set $X=U_c^{-1}\otimes U$ and $Y=V_c^{-1}\otimes V$.
Each dependency now gives positive separated monomials with
weighted product one. Their nonnegative differences are sums of
squares by the identities
\begin{align*}
 p+p^{-1}-2&=(\sqrt p-p^{-1/2})^2,\\
 a^3+b^3+c^3-3abc
 &=\tfrac12(a+b+c)\bigl[(a-b)^2+(b-c)^2+(c-a)^2\bigr],\\
 a^4+b^4+c^4+d^4-4abcd
 &=(a^2-b^2)^2+(c^2-d^2)^2+2(ab-cd)^2.
\end{align*}
Use cube roots for weights $(1,1,1)$, and fourth roots with the
last monomial repeated for weights $(1,1,2)$. Every root is taken
separately in the source and testing fields. Expand each positive
monomial prefactor before taking its square root. Multiplying by
the positive source weights and summing proves the assertion.
\end{proof}

In one field, Lemma~\ref{lem:insight-balance} and
Eq.~\eqref{eq:insight-transverse} show that a nonnegative event
source is physical whenever it has a noncommuting stationary zero:
\begin{equation}\label{eq:insight-contact}
 P_c(u_c,v_c)\psi_c=0,\qquad
 \psi_c^{\mathsf T}P_{c,u}(u_c,v_c)\psi_c=
 \psi_c^{\mathsf T}P_{c,v}(u_c,v_c)\psi_c=0.
\end{equation}
The positive transverse block makes its adjugate a positive
multiple of $\psi_c\psi_c^{\mathsf T}$, so these equations give
the three scalar contact equations used in the lemma.

This also reduces the supporting inequalities. Normalize
$d+\sum_i c_i=1$. Call the angular sum coercive when zero lies in
the interior of the convex hull of its active exponents. Over $\mathbb R$,
it then diverges whenever $(\log U,\log V)$ leaves every bounded set.
A coercive extreme source must attain zero;
otherwise opposite small perturbations remain physical. At a
fixed zero, Eq.~\eqref{eq:insight-contact} is a system of six
homogeneous linear equations in the source coefficients. Extremality
therefore leaves at most six positive event weights. The compatible
omissions in the eight-event cycle leave at least five. Up to the
CHSH symmetries, the resulting patterns are precisely
\begin{equation}\label{eq:insight-two-families}
\begin{aligned}
 \mathsf C_1:&\quad\lambda_{\rm ev}=n=0,\quad
 d,m,C,A,\alpha>0,\quad D,\beta\ge0,\quad D+\beta>0,\\
 \mathsf C_2:&\quad\lambda_{\rm ev}=\beta=0,\quad
 d,m,C,A,D,\alpha>0,\quad n\ge0.
\end{aligned}
\end{equation}
Noncoercive sources belong to two closed families with explicit
small lifts. Chapter~4 proves the extreme-source
reduction, the compatibility assertion and these closed-family
lifts. Chapter~5 proves that the actual contacts have
the required signs and positive native gaps, including $D=0$,
$\beta=0$ and $n=0$. Thus no positive infinitesimal event is
discarded by replacing it with zero.

\section{Recovering the full state polynomial at singular contacts}
\label{sec:insight-normal}

The angular lemma certifies the determinant. It does not yet
certify the full state polynomial in the separated ring. Ordinary
state elimination can introduce an inverse depending on both the
source and the test; such an inverse is not licensed in $T$.
The local argument keeps the state amplitudes in explicit polynomial
identities and controls the coefficients of their Gram matrix.

Let $B_c,B_t$ be the bounded elements of the two fields, and put
$B=B_c\otimes B_t$. The residue $\overline a$ of a bounded tensor
is obtained by taking the real standard part in each factor.

\begin{lemma}[A strict real Gram matrix]\label{lem:insight-gram}
If $G\in\operatorname{Sym}_r(B)$ and $\overline G\succ0$, then
$z^{\mathsf T}Gz\in\Sigma T^2$ for every $z\in T^r$.
The entries of $z$ need not be bounded or independent.
\end{lemma}
\begin{proof}
First suppose $a\in B$ has positive real residue $a_0$.
Write $a=a_0+\sum_jx_j\otimes y_j$, with bounded factors and
an infinitesimal factor in each summand. Distribute $a_0$ into
positive real numbers $a_j$. If $x_j$ is infinitesimal, choose
a real $M_j>|y_j|$ and use
\[
 a_j+x_j\otimes y_j
 ={1\over2}\bigl[(a_j/M_j+x_j)\otimes(M_j+y_j)
                +(a_j/M_j-x_j)\otimes(M_j-y_j)\bigr].
\]
Each pure factor on the right is positive and has a square root
in its own field. Interchange the factors when $y_j$ is
infinitesimal. Thus $a$ is a tensor sum of squares.

A real congruence reduces the matrix statement to
$\overline G=I$. Choose real $\delta>0$ with $(r-1)\delta<1$.
Writing $e_i$ for the coordinate vectors, the exact decomposition
\[
\begin{aligned}
 G={}&\sum_i[G_{ii}-(r-1)\delta]e_ie_i^{\mathsf T}\\
 &+\sum_{i<j}{\delta+G_{ij}\over2}(e_i+e_j)(e_i+e_j)^{\mathsf T}\\
 &+\sum_{i<j}{\delta-G_{ij}\over2}(e_i-e_j)(e_i-e_j)^{\mathsf T}
\end{aligned}
\]
has scalar coefficients of positive residue. Apply the scalar
argument and pair the resulting squares with $z$.
\end{proof}

Every local certificate reduces to this lemma. If $\Pi$ is a sum
of positive event squares and $s=s_0+A\xi$ is the chosen state
substitution, quadratic state degree gives the exact identity
\begin{equation}\label{eq:insight-whole-state}
 \begin{gathered}
 \mathcal R(\xi):=(F-\Pi)(s_0+A\xi)
        =g_0+\sum_i g_i\xi_i+\xi^{\mathsf T}H\xi,\\
 g_0=[1]\mathcal R,\qquad g_i=[\xi_i]\mathcal R,\qquad
 H_{ij}=\tfrac12\partial_{\xi_i}\partial_{\xi_j}\mathcal R.
 \end{gathered}
\end{equation}
The scalar $g_0$, every first state-normal coefficient $g_i$ and
the whole state Hessian $H$ must be included. Monic polynomial
division at the contact expresses their vanishing factors in the
angular roots. The raw monomial support and the source bounds
then show that the resulting coefficient corrections have zero
residue. Contact vanishing alone does not provide either assertion.
The exact restricted division theorem is in
Section~6.2; the common polynomial
coefficient operations are in Section~B.1.1.

A useful uniform step removes a choice between a vanishing state
coefficient and a vanishing angular coefficient.

\begin{lemma}[Complementary source weights]\label{lem:insight-complementary}
Let $a,b\ge0$ be source coefficients with $S=a+b$ a positive
bounded unit. Suppose $\ell,c_i\in B$ have zero residue and
\[
 \sqrt b\,\ell\Gamma=\sum_i c_i z_i.
\]
Then $\ell\Gamma$ is a linear combination, with zero-residue
bounded coefficients, of $\sqrt a\Gamma,z_1,\ldots,z_m$.
This includes either weight equal to zero.
\end{lemma}
\begin{proof}
Use the exact identity
\[
 \ell\Gamma={\sqrt a\,\ell\over S}(\sqrt a\Gamma)
                  +{\sqrt b\over S}\sum_i c_i z_i.
\]
Only the positive source unit $S$ is inverted. Both sets of
coefficients are bounded with zero residue.
\end{proof}

Here is a small polynomial example of the conductor step.
Let $k>0$ be a source infinitesimal, $s>0$ a bounded source
element, and $u,v>0$ independent testing elements. Put
\[
 U_1=u^2v,\quad U_2=kuv^2/s,\quad U_3=k^2s,\quad g=kuv,
 \qquad d_i=U_i-g.
\]
Then $U_1U_2U_3=g^3$. With $p_i=U_i^{1/3}$,
\[
 \sum_i U_i^2-3g^2
 =\sum_i d_i^2+
       g(p_1+p_2+p_3)\sum_{i<j}(p_i-p_j)^2\in\Sigma T^2.
\]
Thus the angular identity supplies three unit diagonal coefficients
on $d_1,d_2,d_3$. The exact polynomial relations
\begin{equation}\label{eq:insight-polynomial-conductors}
 g(u-k)=k d_1,\qquad g(v-s)=s d_2
\end{equation}
express contact factors in those same roots, with no division by
$u-k$ or $v-s$. The first coefficient is infinitesimal. If $s$
is a unit, the second becomes a zero-residue correction only when
the complete state-normal term supplies an additional small
coefficient. That coefficient bound is part of the full state-normal
calculation; stationarity alone does not supply it.

The C1 application, Theorem~B.23,
combines polynomial conductors with complementary source weights.
Its statement gives the exact source domain, and its proof verifies
the complete state-normal bounds uniformly over that domain.
It covers all infinitesimal independent $u,w$, every state ratio,
and both $D=0$ and $\beta=0$; separate complementary results
supply the unit-Alice and other projective tests.

The C2 construction uses the same operations on the full state
polynomial: positive event squares remain whole, polynomial contact
relations produce angular roots, and a strict real matrix controls
the remaining directions. Its shared local and complementary
identities are Sections~C.3.1
and~C.3.2; the bounded-ratio theorem is
Section~C.11.6. The hypotheses of these
statements are proved from the original source coefficients.
They are not consequences of the determinant identity alone.

\begin{proposition}[Complete canonical certificate]
\label{prop:insight-canonical}
Let $K_c,K_t$ be any separately labelled real closed extensions
of $\mathbb R$. For a source in either family
\eqref{eq:insight-two-families} with an actual finite noncommuting
stationary zero~\eqref{eq:insight-contact}, the complete homogeneous parent
belongs to $\Sigma T^2$ at every independent projective testing
pair and every homogeneous real or complex state.
\end{proposition}
\begin{proof}
The full source-dependent proof is in the complete companion.
After normalizing $d=1$, the source ratio
\[
 \delta={d\over d+\lambda_{\rm ev}+n+m+C+A+D+\alpha+\beta}
\]
is either a positive unit or a positive infinitesimal. A unit
gives bounded source coefficients and is covered by
Chapter~A. The infinitesimal case is covered by the
C1 and C2 source theorems in Chapters~B
and~C. Those chapters prove the source partitions
and the coefficient estimates needed in
Eq.~\eqref{eq:insight-whole-state}, and give all complementary
testing and state charts. The noncommuting-residue theorem
(Theorem~6.2) applies at arbitrary source
scale. Literal commuting tests have a positive CHSH slack directly.
The homogeneous parent includes both zero and infinite half-angles.
The proofs retain the permitted literal $D=0$, $\beta=0$ and
$n=0$ faces. Applying the real result separately to the real and
imaginary state parts gives the complex-state result.
\end{proof}

\section{Assembly and the extraction of an exact pencil}
\label{sec:insight-assembly}

The source reduction in Chapter~4 expresses any
valid slack as a nonnegative combination of at most nine event,
closed-family or stationary source slacks. Its fixed finite formula
transfers to $K_c$. Equation~\eqref{eq:insight-finite-model} gives
at most nine testing atoms in $K_t$. Events, the explicit
closed-family lifts and Proposition~\ref{prop:insight-canonical}
certify every source/test atom pair. If
\[
 f=\sum_j a_jf_j,\qquad q=\sum_iw_iq_i,\qquad
 a_j,w_i\ge0,\quad\sum_iw_i=1,\quad
 [f_j,q_i]=\sum_k h_{ijk}^2,
\]
then exactly
\[
 [f,q]=\sum_{i,j,k}
       \bigl[(\sqrt{a_j}\otimes\sqrt{w_i})h_{ijk}\bigr]^2.
\]
The normalized testing weights preserve the affine constant.
All positive source and testing normalizations are
restored by square roots in their own fields, as in Eq.~(RFA13)
of Chapter~2. This proves
Eq.~\eqref{eq:insight-universal-goal}. Scheiderer's criterion now
gives one finite semidefinite representation of all eight coordinates.

Finally, at any fixed $N,p$, equality between the proposed projection
and Eq.~\eqref{eq:insight-finite-model} is a first-order condition
over $\mathbb Q$ on the matrix entries: positive semidefiniteness
is expressed by its principal minors. A real solution therefore
has a solution in the real algebraic field. This proves the
algebraic-coefficient assertion in
Theorem~\ref{thm:insight-finite-lift}. It also gives an exact
extraction procedure in principle: enumerate the finite size pairs,
decide each existence sentence by quantifier elimination, and
return an algebraic sample for the first successful pair.
The theorem ensures termination. No successful full-set pencil or
useful size bound has been evaluated here; finding a small explicit
representation requires further constructive work.
